
\documentclass[12pt,a4paper]{article}
\usepackage{amsmath, amssymb, geometry, graphicx}
\geometry{margin=2cm}

\title{Estudo da Função de Energia Potencial de Lennard-Jones}
\author{}
\date{}

\begin{document}

\maketitle

\section{Definição da Função}

A energia potencial de Lennard-Jones é dada por:

\[
U(r) = 4\varepsilon \left[ \left( \frac{\sigma}{r} \right)^{12} - \left( \frac{\sigma}{r} \right)^6 \right]
\]

onde:
\begin{itemize}
    \item $\varepsilon$: profundidade do poço de potencial;
    \item $\sigma$: distância para a qual $U(r)=0$.
\end{itemize}


\section{Interseção com o eixo x}

\[
U(r) = 0 \quad \Rightarrow \quad r = \sigma
\]


\section{Mínimo da Função}

Derivando:

\[
\frac{dU}{dr} = 4\varepsilon \left( -12 \frac{\sigma^{12}}{r^{13}} + 6 \frac{\sigma^6}{r^7} \right)
\]

Impondo $\frac{dU}{dr}=0$:

\[
r = \sigma \cdot 2^{1/6} \approx 1,122\sigma
\]

Neste ponto:

\[
U_{\min} = -\varepsilon
\]


\section{Ponto de Inflerão}

A segunda derivada é:

\[
\frac{d^2U}{dr^2} = 4\varepsilon \left( 156 \frac{\sigma^{12}}{r^{14}} - 42 \frac{\sigma^6}{r^8} \right)
\]

Igualando a zero:

\[
r = \sigma \left( \frac{26}{7} \right)^{1/6} \approx 1,244\sigma
\]


\section{Comportamento Assintótico}

\[
\lim_{r \to 0^+} U(r) = +\infty
\qquad
\lim_{r \to \infty} U(r) = 0^-
\]


\section{Força}
\begin{flushleft}
A força é dada por:

\[
F(r) = -\frac{dU}{dr} = 24\varepsilon \left( \frac{2\sigma^{12}}{r^{13}} - \frac{\sigma^6}{r^7} \right)
\]

- Atrativa ($F(r)<0$) para $r > 1,122\sigma$\\  
- Repulsiva ($F(r)>0$) para $r < 1,122\sigma$\\  
\end{flushleft}

\section{Termos Individuais}

\[
U_{\text{rep}}(r) = 4\varepsilon \left( \frac{\sigma}{r} \right)^{12}
\qquad
U_{\text{atr}}(r) = -4\varepsilon \left( \frac{\sigma}{r} \right)^6
\]

- Para $r \ll \sigma$: termo repulsivo domina;  
- Para $r \gg \sigma$: termo atrativo domina;  


\end{document}
